\chapter{Completion Queues and Non-Lossy Results}
\label{chap:completion-queues}

A completion queue (CQ) carries asynchronous kernel-operation results to
userspace. Its central safety rule is that a full userspace ring must not silently
discard a terminal result.

\section{When to use a completion queue}

Use submission/completion when the operation belongs to the \textbf{kernel or a
hardware-facing mechanism}, not to another userspace service:

\begin{itemize}
    \item Timer expiry.
    \item DMA completion (disk read finished, packet transmitted).
    \item IRQ-derived device completion (data available on UART).
    \item Asynchronous page operations.
    \item Waiting for thread or process exit.
    \item Kernel-managed asynchronous object work.
\end{itemize}

Use endpoint IPC (\autoref{chap:endpoints}) when one protection domain invokes another.
Confusing these two paths leads to architectural mistakes: a service call is not a
kernel operation, and a reply token is not a completion record.

\section{The CQ ring --- a shared-memory completion buffer}

A completion queue is a \textbf{shared-memory ring buffer} between the kernel
(producer) and userspace (consumer). A single 4 KiB page contains a header and a
circular array of fixed-size entries.

Each entry is 32 bytes:

\begin{lstlisting}[style=rustCode]
struct CompletionEntry {
    operation: u64,    // operation identifier (e.g., timer id)
    cookie: u64,       // opaque value set by the submitter
    status: u32,       // completion status code
    flags: u32,        // per-entry flags
    result: i64,       // operation result (e.g., bytes read)
}
\end{lstlisting}

The ring page is mapped read-write in userspace because the consumer advances its
tail; the kernel accesses the same frame through its privileged mapping. The ABI
assigns each side its own fields:

\begin{itemize}
    \item Userspace drains entries without syscalls --- advancing the
        consumer tail is a store to a shared page; reading entries is a load.
        No kernel crossing needed for the common fast path.
    \item The kernel publishes entries with a release fence --- the consumer
        sees coherent data without additional synchronization.
    \item Overflow is detected, not hidden --- an overflow counter in the
        header lets userspace detect when entries were produced faster than
        consumed.
\end{itemize}

\section{The non-lossy invariant}

\begin{quote}
\textbf{A terminal operation result must never be silently lost because the
userspace CQ ring is full.}
\end{quote}

This is an architectural invariant, so when the ring is full, the
kernel has several options:

\begin{itemize}
    \item Retain overflow in a per-domain kernel backlog --- entries that
        cannot fit in the ring wait in a kernel-side buffer. When userspace
        drains the ring, the backlog drains into it.
    \item Stop accepting new submissions --- the kernel may return
        \texttt{WouldBlock} on \texttt{submit}, applying backpressure at the
        submission point rather than at the completion point.
    \item Propagate \texttt{WouldBlock} --- the CQ's wait operation can
        indicate that the ring is full and must be drained.
\end{itemize}

The overflow counter is pressure telemetry: it tells the consumer that it is
falling behind.

In a critical system, a lost completion for a disk write or a timer
expiry is a correctness bug with potentially severe consequences. CharlotteOS
treats each completion as a datum that must be delivered at least once.

\section{Per-shard CQ partitioning}

Each address space owns a set of completion queues, keyed by \rustfn{CqId}. Queue
0 is the default. For a service with multiple sitas shards, \textbf{per-shard CQ
rings} enable targeted wakes:

\begin{itemize}
    \item Shard A's executor waits on CQ 0.
    \item Shard B's executor waits on CQ 1.
    \item A completion destined for Shard A is posted to CQ 0.
    \item Only Shard A's executor is woken; Shard B is undisturbed.
\end{itemize}

This prevents cross-shard completion stealing (a pathological behavior where one
busy shard drains entries intended for another) and keeps wakeup latency
proportional to the number of entries, not the number of shards.

\section{Operation lifecycle}

Every asynchronous operation transitions through a defined state machine:

\begin{verbatim}
Created
    -> Submitted
    -> Accepted | Rejected
    -> InFlight
    -> Completed | Failed | Cancelled
    -> ResultObserved
    -> ResourcesReclaimed
\end{verbatim}

Each transition is explicit and observable:

\begin{description}
    \item[Created] The operation descriptor exists in userspace but has not been
        submitted to the kernel.
    \item[Submitted] \texttt{submit()} has been called. The kernel has the
        operation record but has not yet validated it.
    \item[Accepted / Rejected] The kernel has validated the operation. It may be
        rejected immediately (invalid parameters, insufficient rights, queue
        full) with a synchronous error.
    \item[InFlight] The operation is being processed. The kernel or device owns
        any associated buffers.
    \item[Completed / Failed / Cancelled] A terminal state. The result is posted
        to the CQ. Buffer ownership returns to userspace.
    \item[ResultObserved] The consumer has read the CQ entry.
    \item[ResourcesReclaimed] The operation's capability-table slot has been
        freed (via \texttt{close}).
\end{description}

\section{Waiting on a CQ --- \texttt{CQ\_WAIT}}

When a shard's executor has no ready tasks, it calls \texttt{CQ\_WAIT} on its
assigned queue. The kernel blocks the calling thread until one of three things
happens:

\begin{enumerate}
    \item A completion entry is posted to the CQ.
    \item An explicit \texttt{CQ\_WAKE} arrives (from another shard or the
        kernel).
    \item A deadline elapses (timeout variant: \texttt{CQ\_WAIT\_TIMEOUT}).
\end{enumerate}

The wait is race-free. The sequence is:

\begin{enumerate}
    \item Inspect the CQ ring --- if completions are pending, return immediately.
    \item Register a waiter against the CQ's generation counter.
    \item Re-inspect the CQ ring --- if completions arrived during registration,
        unregister and return.
    \item Block the thread --- it will be woken when the generation counter changes
        (a new entry is posted or a \texttt{CQ\_WAKE} bumps it).
\end{enumerate}

This is the \textbf{single wait point} for the entire shard. Three event sources
converge on it:

\begin{itemize}
    \item Kernel completions (CQ entries from \texttt{submit} operations).
    \item Endpoint readiness (a coalesced wake when messages arrive on a bound
        endpoint).
    \item Device interrupts (a coalesced wake from a bound interrupt object).
\end{itemize}

The shard does not need separate wait operations for timers, IPC messages, and
device events. One aggregated wait covers all of them.

\section{LP affinity for CQ waits}

A thread that blocks in \texttt{CQ\_WAIT} is assigned an affinity LP. When woken,
the scheduler returns it to the same LP rather than scanning for the least-loaded
one. This has three benefits:

\begin{enumerate}
    \item Timer events stay on the same LP's queue, eliminating
        cross-LP migration races (observer-not-found when a completion is posted
        on a different LP than the one where the waiter was registered).
    \item Cache warmth is preserved --- the waking thread finds its data
        in the LP's local caches.
    \item Shard-to-LP affinity is maintained --- the Sitas shard that blocked
        on CQ 0 on LP 2 wakes up on LP 2, where its shard-local state lives.
\end{enumerate}

\section{Operation IDs versus per-operation capabilities}

An earlier experimental design returned a \texttt{CompletionCap} for every
submission. The final ABI uses \textbf{operation IDs} for the common path:

\begin{itemize}
    \item \texttt{submit()} returns an \rustfn{OperationId} (a lightweight
        identifier).
    \item The result arrives in the CQ with the same operation ID.
    \item No capability-table slot is consumed for common operations.
    \item A first-class operation capability is created only when independent
        waiting, cancellation, delegation, inspection, or policy requires it.
\end{itemize}

This avoids capability-table pressure for high-rate operations (timer events,
network packet completions) while preserving the capability model for operations
that need first-class authority.

\section{Cancellation and buffer ownership}

When a userspace task drops a completion handle or explicitly cancels an
in-flight operation:

\begin{enumerate}
    \item The kernel marks the operation as \texttt{Cancelling} but \textbf{retains
        any transferred buffer.} The buffer may still be the target of DMA or
        kernel access.
    \item A terminal \texttt{Cancelled} completion is eventually posted to the CQ,
        returning buffer ownership to userspace.
    \item On domain teardown (process exit), outstanding buffers may be leaked
        rather than freed while the kernel or a device may still touch them.
        Reclamation goes through explicit operation and device teardown.
\end{enumerate}

Under this contract, a buffer is retained while the kernel or device may still
reach it. The safe Rust wrapper consumes the owned value on submission and returns
it on terminal completion.

\section{Why completion queues matter for critical software}

\begin{enumerate}
    \item Non-lossy completions are a correctness requirement --- a disk-write
        confirmation, a timer expiry, or a sensor reading is a fact and losing it
        is a bug. The backlog retains terminal results when the ring is full.

    \item The single aggregated wait eliminates a class of race conditions ---
        a shard that must wait on ``timer or IPC message or device interrupt'' in
        a conventional OS must manage three separate wait objects, with potential
        lost-wake races between checking one and arming another. One CQ means one
        atomic check-and-arm sequence.

    \item Per-shard CQ rings enable predictable latency --- no shard wakes
        another shard's executor to deliver a completion. Wakes are targeted,
        bounded, and proportional to the actual work.

    \item The CQ ring is zero-syscall in the common case --- the fast path
        --- drain entries, process results, loop --- requires no kernel
        transitions. Only when the ring is empty and the shard must block does a
        syscall occur.

    \item Cancellation is defined, not best-effort --- every cancellation
        path ends in a terminal completion. There is no ``the buffer might or
        might not be freed'' ambiguity. The ownership contract is explicit and
        testable.
\end{enumerate}

\endinput
